% GENERATED FILE. Edit canonical lessons, then run npm run generate:books.
% canonical-source-hash: fnv1a64:c298d3d41dabab62
% canonical-lessons: HI-C451-sair, HI-C451-paryatak, HI-C451-nazara, HI-C451-smarak, HI-C451-niharna, HI-R451-review-445-451, HI-R451-review-final-451

\chapter{Outings and Views}
\label{ch:hi-outings-and-views}
\hlchaptermodality{451}

\begin{chapteropening}
\textbf{By the end of this chapter:} \emph{I can talk about an outing, tourists, a view and a monument, and gaze at the hills.}

Complete the last lesson of chapter 451: I can talk about an outing, tourists, a view and a monument, and gaze at the hills.
\end{chapteropening}

\section[\texorpdfstring{\hi{सैर} (sair)}{सैर (sair)}]{\hi{सैर} (sair) — an outing; a stroll}

\label{lesson:HI-C451-sair}



\begin{quote}
Before the new one: say the Hindi for to grow (of crops), then the Hindi for to milk.
\end{quote}

\subsection*{You'll want to know: \hi{सैर}}
\textbf{\hi{सैर}} — \emph{sair} — \textquotedblleft{}an outing; a stroll\textquotedblright{}.

\begin{quote}
\hi{छुट्टी} \hi{में} \hi{हम} \hi{पहाड़ों} \hi{की} \hi{सैर} \hi{पर} \hi{गए।}
\end{quote}

\begin{quote}
\emph{chuṭṭī meṁ ham pahāṛoṁ kī sair par gae.}
\end{quote}

\begin{quote}
On holiday we went on a trip to the hills.
\end{quote}

Feminine.

\begin{cousinweb}
\textbf{\hi{सैर}} (\emph{sair}) is usually traced to Arabic \emph{sayr}, \textquotedblleft{}travel, going about\textquotedblright{}, through Persian: a trip for pleasure.
\end{cousinweb}

\begin{grammarlens}[title={What you've built}]
One more word for getting around.
\end{grammarlens}

\subsection*{Your turn}
\begin{itemize}
\raggedright
  \item \emph{Say it:} \emph{sair}
  \item \emph{Say it:} \emph{sair}, once more
  \item \emph{Say it:} \emph{chuṭṭī meṁ ham pahāṛoṁ kī sair par gae.} — On holiday we went on a trip to the hills.
  \item \emph{Recall:} say the Hindi for to grow (of crops), then the Hindi for to milk, then say \emph{sair} again
\end{itemize}

\begin{tcolorbox}[breakable,colback=teal!4,colframe=teal!35!black,title={Before you move on}]
What does \hi{सैर} mean? (\textquotedblleft{}An outing; a stroll\textquotedblright{}.) Say it once more.
\end{tcolorbox}
\section[\texorpdfstring{\hi{पर्यटक} (paryaṭak)}{पर्यटक (paryaṭak)}]{\hi{पर्यटक} (paryaṭak) — a tourist}

\label{lesson:HI-C451-paryatak}



\begin{quote}
Before the new one: say the Hindi for to milk, then the Hindi for an outing.
\end{quote}

\subsection*{You'll want to know: \hi{पर्यटक}}
\textbf{\hi{पर्यटक}} — \emph{paryaṭak} — \textquotedblleft{}a tourist\textquotedblright{}.

\begin{quote}
\hi{गर्मियों} \hi{में} \hi{यहाँ} \hi{बहुत} \hi{पर्यटक} \hi{आते} \hi{हैं।}
\end{quote}

\begin{quote}
\emph{garmiyoṁ meṁ yahā̃ bahut paryaṭak āte haiṁ.}
\end{quote}

\begin{quote}
A lot of tourists come here in the summer.
\end{quote}

Masculine or feminine, by the person meant.

\begin{cousinweb}
\textbf{\hi{पर्यटक}} (\emph{paryaṭak}) is a tatsama, Sanskrit \emph{paryaṭaka}, one who wanders about (Monier-Williams gives a tramp, a vagabond); the modern sense of a tourist is the newer one.
\end{cousinweb}

\begin{grammarlens}[title={What you've built}]
One more word for getting around.
\end{grammarlens}

\subsection*{Your turn}
\begin{itemize}
\raggedright
  \item \emph{Say it:} \emph{paryaṭak}
  \item \emph{Say it:} \emph{paryaṭak}, once more
  \item \emph{Say it:} \emph{garmiyoṁ meṁ yahā̃ bahut paryaṭak āte haiṁ.} — A lot of tourists come here in the summer.
  \item \emph{Recall:} say the Hindi for to milk, then the Hindi for an outing, then say \emph{paryaṭak} again
\end{itemize}

\begin{tcolorbox}[breakable,colback=teal!4,colframe=teal!35!black,title={Before you move on}]
What does \hi{पर्यटक} mean? (\textquotedblleft{}A tourist\textquotedblright{}.) Say it once more.
\end{tcolorbox}
\section[\texorpdfstring{\hi{नज़ारा} (nazārā)}{नज़ारा (nazārā)}]{\hi{नज़ारा} (nazārā) — a view, a sight}

\label{lesson:HI-C451-nazara}



\begin{quote}
Before the new one: say the Hindi for an outing, then the Hindi for a tourist.
\end{quote}

\subsection*{You'll want to know: \hi{नज़ारा}}
\textbf{\hi{नज़ारा}} — \emph{nazārā} — \textquotedblleft{}a view, a sight\textquotedblright{}.

\begin{quote}
\hi{पहाड़} \hi{से} \hi{शहर} \hi{का} \hi{नज़ारा} \hi{बहुत} \hi{सुंदर} \hi{था।}
\end{quote}

\begin{quote}
\emph{pahāṛ se śahar kā nazārā bahut sundar thā.}
\end{quote}

\begin{quote}
The view of the city from the mountain was beautiful.
\end{quote}

Masculine.

\begin{cousinweb}
\textbf{\hi{नज़ारा}} (\emph{nazārā}) is usually traced to Arabic \emph{naẓāra}, \textquotedblleft{}a looking, a sight\textquotedblright{}, through Persian; the root is the same as in \hi{नज़र}, a glance.
\end{cousinweb}

\begin{grammarlens}[title={What you've built}]
One more word for getting around.
\end{grammarlens}

\subsection*{Your turn}
\begin{itemize}
\raggedright
  \item \emph{Say it:} \emph{nazārā}
  \item \emph{Say it:} \emph{nazārā}, once more
  \item \emph{Say it:} \emph{pahāṛ se śahar kā nazārā bahut sundar thā.} — The view of the city from the mountain was beautiful.
  \item \emph{Recall:} say the Hindi for an outing, then the Hindi for a tourist, then say \emph{nazārā} again
\end{itemize}

\begin{tcolorbox}[breakable,colback=teal!4,colframe=teal!35!black,title={Before you move on}]
What does \hi{नज़ारा} mean? (\textquotedblleft{}A view, a sight\textquotedblright{}.) Say it once more.
\end{tcolorbox}
\section[\texorpdfstring{\hi{स्मारक} (smārak)}{स्मारक (smārak)}]{\hi{स्मारक} (smārak) — a monument, a memorial}

\label{lesson:HI-C451-smarak}



\begin{quote}
Before the new one: say the Hindi for a tourist, then the Hindi for a view.
\end{quote}

\subsection*{You'll want to know: \hi{स्मारक}}
\textbf{\hi{स्मारक}} — \emph{smārak} — \textquotedblleft{}a monument, a memorial\textquotedblright{}.

\begin{quote}
\hi{शहर} \hi{के} \hi{बीच} \hi{में} \hi{एक} \hi{पुराना} \hi{स्मारक} \hi{है।}
\end{quote}

\begin{quote}
\emph{śahar ke bīc meṁ ek purānā smārak hai.}
\end{quote}

\begin{quote}
There is an old monument in the middle of the town.
\end{quote}

Masculine.

\begin{cousinweb}
\textbf{\hi{स्मारक}} (\emph{smārak}) is a tatsama, Sanskrit \emph{smāraka}, \textquotedblleft{}recalling, reminding of\textquotedblright{} (Monier-Williams), from the root \emph{smṛ}, \textquotedblleft{}to remember\textquotedblright{}: a monument is what reminds.
\end{cousinweb}

\begin{grammarlens}[title={What you've built}]
One more word for getting around.
\end{grammarlens}

\subsection*{Your turn}
\begin{itemize}
\raggedright
  \item \emph{Say it:} \emph{smārak}
  \item \emph{Say it:} \emph{smārak}, once more
  \item \emph{Say it:} \emph{śahar ke bīc meṁ ek purānā smārak hai.} — There is an old monument in the middle of the town.
  \item \emph{Recall:} say the Hindi for a tourist, then the Hindi for a view, then say \emph{smārak} again
\end{itemize}

\begin{tcolorbox}[breakable,colback=teal!4,colframe=teal!35!black,title={Before you move on}]
What does \hi{स्मारक} mean? (\textquotedblleft{}A monument, a memorial\textquotedblright{}.) Say it once more.
\end{tcolorbox}
\section[\texorpdfstring{\hi{निहारना} (nihārnā)}{निहारना (nihārnā)}]{\hi{निहारना} (nihārnā) — to gaze at (admiringly)}

\label{lesson:HI-C451-niharna}



\begin{quote}
Before the new one: say the Hindi for a view, then the Hindi for a monument.
\end{quote}

\subsection*{You'll want to know: \hi{निहारना}}
\textbf{\hi{निहारना}} — \emph{nihārnā} — \textquotedblleft{}to gaze at (admiringly)\textquotedblright{}.

\begin{quote}
\hi{मीरा} \hi{ने} \hi{देर} \hi{तक} \hi{पहाड़ों} \hi{को} \hi{निहारा।}
\end{quote}

\begin{quote}
\emph{Mīrā ne der tak pahāṛoṁ ko nihārā.}
\end{quote}

\begin{quote}
Mira gazed at the mountains for a long time.
\end{quote}

Past, with \hi{ने} (he or she …): \hi{उसने} \hi{सूर्योदय} \hi{को} \hi{निहारा}, \emph{usne sūryoday ko nihārā}.

\begin{cousinweb}
\textbf{\hi{निहारना}} (\emph{nihārnā}) is a tadbhava of Sanskrit \emph{nibhālayati}, \textquotedblleft{}perceives\textquotedblright{} (CDIAL 7228), which lists Hindi \emph{nihārnā}, \textquotedblleft{}to look at\textquotedblright{}: to look long, and with pleasure.
\end{cousinweb}

\begin{grammarlens}[title={What you've built}]
One more word for getting around.
\end{grammarlens}

\subsection*{Your turn}
\begin{itemize}
\raggedright
  \item \emph{Say it:} \emph{nihārnā}
  \item \emph{Say it:} \emph{nihārnā}, once more
  \item \emph{Say it:} \emph{Mīrā ne der tak pahāṛoṁ ko nihārā.} — Mira gazed at the mountains for a long time.
  \item \emph{Recall:} say the Hindi for a view, then the Hindi for a monument, then say \emph{nihārnā} again
\end{itemize}

\begin{tcolorbox}[breakable,colback=teal!4,colframe=teal!35!black,title={Before you move on}]
What does \hi{निहारना} mean? (\textquotedblleft{}To gaze at (admiringly)\textquotedblright{}.) Say it once more.
\end{tcolorbox}
\section[(dialogue)]{Looking back: from Boats, Nets and Shells to Outings and Views}

\label{lesson:HI-R451-review-445-451}



\begin{quote}
Say the Hindi for a small boat and a boatman.
\end{quote}

\subsection*{Your turn}
Say these aloud:

\begin{itemize}
\raggedright
  \item a small boat, a boatman, a net, the high tide, a shell
  \item nature, the environment, a plant, to bloom, to wilt
  \item a bower, a creeper, moss, an insect, to wind round
  \item a beak, a wing, hunting, to peck, to prick
  \item a jackal, a firefly, a cricket, to graze, to take to graze
  \item to be shed, to be uprooted, to plough, to grow (of crops), to milk
  \item an outing, a tourist, a view, a monument, to gaze at
\end{itemize}

\begin{tcolorbox}[breakable,colback=teal!4,colframe=teal!35!black,title={Before you move on}]
A small boat? (\textbf{\hi{किश्ती}}, \emph{kiśtī}.) A boatman? (\textbf{\hi{मल्लाह}}, \emph{mallāh}.) A net? (\textbf{\hi{जाल}}, \emph{jāl}.) The high tide? (\textbf{\hi{ज्वार}}, \emph{jvār}.) A shell? (\textbf{\hi{सीप}}, \emph{sīp}.) Nature? (\textbf{\hi{प्रकृति}}, \emph{prakṛti}.) The environment? (\textbf{\hi{पर्यावरण}}, \emph{paryāvaraṇ}.) A plant? (\textbf{\hi{पौधा}}, \emph{paudhā}.) To bloom? (\textbf{\hi{खिलना}}, \emph{khilnā}.) To wilt? (\textbf{\hi{कुम्हलाना}}, \emph{kumhlānā}.) A bower? (\textbf{\hi{कुंज}}, \emph{kuñj}.) A creeper? (\textbf{\hi{बेल}}, \emph{bel}.) Moss? (\textbf{\hi{काई}}, \emph{kāī}.) An insect? (\textbf{\hi{कीड़ा}}, \emph{kīṛā}.) To wind round? (\textbf{\hi{लिपटना}}, \emph{lipaṭnā}.) A beak? (\textbf{\hi{चोंच}}, \emph{coñc}.) A wing? (\textbf{\hi{पंख}}, \emph{paṅkh}.) Hunting? (\textbf{\hi{शिकार}}, \emph{śikār}.) To peck? (\textbf{\hi{चुगना}}, \emph{cugnā}.) To prick? (\textbf{\hi{चुभना}}, \emph{cubhnā}.) A jackal? (\textbf{\hi{सियार}}, \emph{siyār}.) A firefly? (\textbf{\hi{जुगनू}}, \emph{jugnū}.) A cricket? (\textbf{\hi{झींगुर}}, \emph{jhī̃gur}.) To graze? (\textbf{\hi{चरना}}, \emph{carnā}.) To take to graze? (\textbf{\hi{चराना}}, \emph{carānā}.) To be shed? (\textbf{\hi{झड़ना}}, \emph{jhaṛnā}.) To be uprooted? (\textbf{\hi{उखड़ना}}, \emph{ukhaṛnā}.) To plough? (\textbf{\hi{जोतना}}, \emph{jotnā}.) To grow (of crops)? (\textbf{\hi{उपजना}}, \emph{upajnā}.) To milk? (\textbf{\hi{दुहना}}, \emph{duhnā}.) An outing? (\textbf{\hi{सैर}}, \emph{sair}.) A tourist? (\textbf{\hi{पर्यटक}}, \emph{paryaṭak}.) A view? (\textbf{\hi{नज़ारा}}, \emph{nazārā}.) A monument? (\textbf{\hi{स्मारक}}, \emph{smārak}.) To gaze at? (\textbf{\hi{निहारना}}, \emph{nihārnā}.)
\end{tcolorbox}
\section[(dialogue)]{Outings and Views — once more}

\label{lesson:HI-R451-review-final-451}



\begin{quote}
Say the Hindi for an outing and a tourist.
\end{quote}

\subsection*{Your turn}
\begin{itemize}
\raggedright
  \item \emph{Say it:} an outing, a tourist, a view, a monument, to gaze at
\end{itemize}

\begin{tcolorbox}[breakable,colback=teal!4,colframe=teal!35!black,title={Before you move on}]
An outing? (\textbf{\hi{सैर}}, \emph{sair}.) A tourist? (\textbf{\hi{पर्यटक}}, \emph{paryaṭak}.) A view? (\textbf{\hi{नज़ारा}}, \emph{nazārā}.) A monument? (\textbf{\hi{स्मारक}}, \emph{smārak}.) To gaze at? (\textbf{\hi{निहारना}}, \emph{nihārnā}.)
\end{tcolorbox}
